% =====================================================================
%  Appendix G: Computational Labs Quick Reference
% =====================================================================

\chapter{Computational Labs Quick Reference}
\label{app:computational-labs}

Computers are useful in multivariable calculus for three jobs:

\[
\begin{array}{c|c}
\text{Job} & \text{Use} \\ \hline
\text{visualize} & \text{surfaces, contours, vector fields, flows}\\[1mm]
\text{approximate} & \text{derivatives, integrals, differential equations}\\[1mm]
\text{check algebra} & \text{symbolic derivatives, Jacobians, line and surface integrals}
\end{array}
\]

This appendix gives compact lab recipes.  It uses Python-style notation, with
\texttt{NumPy}, \texttt{Matplotlib}, \texttt{SciPy}, and \texttt{SymPy}.  The mathematics
does not depend on the software.

\[
\begin{array}{c|c}
\text{Mathematical task} & \text{Computational tool} \\ \hline
z=f(x,y) & \text{surface plot}\\[1mm]
f(x,y)=c & \text{contour plot}\\[1mm]
\mathbf F(x,y) & \text{quiver or stream plot}\\[1mm]
DF(\mathbf p) & \text{finite differences or symbolic Jacobian}\\[1mm]
\iint_R f\,dA & \text{iterated integral, midpoint rule, Monte Carlo}\\[1mm]
\int_C \omega & \text{parametrize and substitute}\\[1mm]
\iint_S \Phi & \text{parametrize and use }\mathbf r_u\times\mathbf r_v\\[1mm]
\mathbf r'=\mathbf F(\mathbf r) & \text{ODE solver}
\end{array}
\]

% ---------------------------------------------------------------------
\section{Plotting surfaces and contour maps}
\label{app:G.1}

\subsection*{Standard imports}

\begin{verbatim}
import numpy as np
import matplotlib.pyplot as plt
\end{verbatim}

For symbolic work:
\[
\boxed{
    \texttt{import sympy as sp}
}
\]

For numerical integration and ODEs:
\[
\boxed{
    \texttt{from scipy.integrate import dblquad, tplquad, solve\_ivp}
}
\]

\subsection*{Grid for a function of two variables}

For
\[
    z=f(x,y),
\]
make arrays of \(x\)- and \(y\)-values:

\begin{verbatim}
x = np.linspace(a, b, 200)
y = np.linspace(c, d, 200)
X, Y = np.meshgrid(x, y)
Z = f(X, Y)
\end{verbatim}

Here \(X\), \(Y\), and \(Z\) are arrays.  The function \(f\) must accept arrays.

\subsection*{Surface plot}

\begin{verbatim}
fig = plt.figure()
ax = fig.add_subplot(projection="3d")

ax.plot_surface(X, Y, Z)

ax.set_xlabel("x")
ax.set_ylabel("y")
ax.set_zlabel("z")
plt.show()
\end{verbatim}

\subsection*{Wireframe plot}

\begin{verbatim}
fig = plt.figure()
ax = fig.add_subplot(projection="3d")

ax.plot_wireframe(X, Y, Z)

plt.show()
\end{verbatim}

A wireframe is often easier to read than a solid surface.

\subsection*{Contour plot}

A contour plot shows level curves:
\[
\boxed{
    f(x,y)=c.
}
\]

\begin{verbatim}
plt.contour(X, Y, Z, levels=20)
plt.xlabel("x")
plt.ylabel("y")
plt.axis("equal")
plt.show()
\end{verbatim}

Filled contour plot:
\[
\boxed{
    \texttt{contourf}
}
\]

\begin{verbatim}
plt.contourf(X, Y, Z, levels=30)
plt.colorbar()
plt.axis("equal")
plt.show()
\end{verbatim}

\subsection*{Implicit curves}

To plot
\[
    F(x,y)=0,
\]
plot the zero contour:

\begin{verbatim}
Z = F(X, Y)
plt.contour(X, Y, Z, levels=[0])
plt.axis("equal")
plt.show()
\end{verbatim}

To plot
\[
    F(x,y)=c,
\]
use:

\begin{verbatim}
plt.contour(X, Y, Z, levels=[c])
\end{verbatim}

\subsection*{Common examples}

Plane:
\[
\boxed{
    z=ax+by+c.
}
\]

Bowl:
\[
\boxed{
    z=x^2+y^2.
}
\]

Saddle:
\[
\boxed{
    z=x^2-y^2.
}
\]

Hill:
\[
\boxed{
    z=e^{-(x^2+y^2)}.
}
\]

Radial wave:
\[
\boxed{
    z=\sin\sqrt{x^2+y^2}.
}
\]

\subsection*{Plotting a parametrized surface}

For
\[
    \mathbf r(u,v)=(x(u,v),y(u,v),z(u,v)),
\]
use a parameter grid:

\begin{verbatim}
u = np.linspace(u0, u1, 150)
v = np.linspace(v0, v1, 150)
U, V = np.meshgrid(u, v)

X = x(U, V)
Y = y(U, V)
Z = z(U, V)

fig = plt.figure()
ax = fig.add_subplot(projection="3d")
ax.plot_surface(X, Y, Z)
plt.show()
\end{verbatim}

\subsection*{Cylinder}

\[
    \mathbf r(\theta,z)=(R\cos\theta,R\sin\theta,z).
\]

\begin{verbatim}
theta = np.linspace(0, 2*np.pi, 150)
z = np.linspace(z0, z1, 100)
Theta, Z = np.meshgrid(theta, z)

X = R*np.cos(Theta)
Y = R*np.sin(Theta)

ax.plot_surface(X, Y, Z)
\end{verbatim}

\subsection*{Sphere}

\[
    \mathbf r(\phi,\theta)
    =
    (R\sin\phi\cos\theta,\ R\sin\phi\sin\theta,\ R\cos\phi).
\]

\begin{verbatim}
phi = np.linspace(0, np.pi, 150)
theta = np.linspace(0, 2*np.pi, 150)
Phi, Theta = np.meshgrid(phi, theta)

X = R*np.sin(Phi)*np.cos(Theta)
Y = R*np.sin(Phi)*np.sin(Theta)
Z = R*np.cos(Phi)

ax.plot_surface(X, Y, Z)
\end{verbatim}

\subsection*{Plotting checklist}

\[
\begin{array}{ll}
1. & \text{Choose a domain large enough to show the important feature.}\\[1mm]
2. & \text{Use enough grid points to avoid jagged artifacts.}\\[1mm]
3. & \text{Use }\texttt{axis("equal")}\text{ for planar geometry.}\\[1mm]
4. & \text{Label axes.}\\[1mm]
5. & \text{Check whether the plot hides singularities or holes.}
\end{array}
\]

\subsection*{Common warnings}

A computer plot is a sample, not a proof.

A surface may look smooth because the grid is too coarse.

A contour may disappear if the requested level is outside the range of plotted values.

A \(3\)-D plot may distort scale unless the axes are adjusted.

% ---------------------------------------------------------------------
\section{Visualizing vector fields}
\label{app:G.2}

\subsection*{Vector field in the plane}

A vector field in the plane has the form
\[
\boxed{
    \mathbf F(x,y)=(P(x,y),Q(x,y)).
}
\]

On a grid:

\begin{verbatim}
x = np.linspace(a, b, 25)
y = np.linspace(c, d, 25)
X, Y = np.meshgrid(x, y)

P = Pfun(X, Y)
Q = Qfun(X, Y)
\end{verbatim}

\subsection*{Quiver plot}

\begin{verbatim}
plt.quiver(X, Y, P, Q)
plt.axis("equal")
plt.xlabel("x")
plt.ylabel("y")
plt.show()
\end{verbatim}

\subsection*{Normalized arrows}

Sometimes large arrows hide the pattern.  Normalize arrows:

\begin{verbatim}
speed = np.sqrt(P**2 + Q**2)
P0 = P / (speed + 1e-12)
Q0 = Q / (speed + 1e-12)

plt.quiver(X, Y, P0, Q0)
plt.axis("equal")
plt.show()
\end{verbatim}

This shows direction but not magnitude.

\subsection*{Stream plot}

Streamlines solve
\[
\boxed{
    \mathbf r'(t)=\mathbf F(\mathbf r(t)).
}
\]

\begin{verbatim}
plt.streamplot(X, Y, P, Q)
plt.axis("equal")
plt.show()
\end{verbatim}

\subsection*{Common vector fields}

Constant field:
\[
\boxed{
    \mathbf F(x,y)=(1,0).
}
\]

Radial source:
\[
\boxed{
    \mathbf F(x,y)=(x,y).
}
\]

Radial sink:
\[
\boxed{
    \mathbf F(x,y)=(-x,-y).
}
\]

Rotation:
\[
\boxed{
    \mathbf F(x,y)=(-y,x).
}
\]

Saddle field:
\[
\boxed{
    \mathbf F(x,y)=(x,-y).
}
\]

Inverse radial field:
\[
\boxed{
    \mathbf F(x,y)=\frac{(x,y)}{x^2+y^2}.
}
\]

\subsection*{Gradient field}

For
\[
    f=f(x,y),
\]
\[
\boxed{
    \nabla f=(f_x,f_y).
}
\]

Numerically, use \texttt{np.gradient} if \(f\) is sampled on a grid:

\begin{verbatim}
Z = f(X, Y)

dy = y[1] - y[0]
dx = x[1] - x[0]

Z_y, Z_x = np.gradient(Z, dy, dx)

plt.quiver(X, Y, Z_x, Z_y)
plt.contour(X, Y, Z, levels=20)
plt.axis("equal")
plt.show()
\end{verbatim}

The order is important:
\[
\boxed{
    \texttt{np.gradient(Z, dy, dx)}
}
\]
returns derivatives in the array-axis order, usually \(y\) first, then \(x\).

\subsection*{Vector field in space}

A vector field in space has the form
\[
\boxed{
    \mathbf F(x,y,z)=(P,Q,R).
}
\]

A \(3\)-D quiver plot:

\begin{verbatim}
fig = plt.figure()
ax = fig.add_subplot(projection="3d")

ax.quiver(X, Y, Z, P, Q, R)

plt.show()
\end{verbatim}

Use fewer grid points in \(3\)-D.  Too many arrows quickly become unreadable.

\subsection*{Field line through one point}

For
\[
    \mathbf r'(t)=\mathbf F(\mathbf r(t)),
\]
use an ODE solver:

\begin{verbatim}
from scipy.integrate import solve_ivp

def ode(t, r):
    x, y = r
    return [Pfun(x, y), Qfun(x, y)]

sol = solve_ivp(ode, [t0, t1], [x0, y0], max_step=0.02)

plt.plot(sol.y[0], sol.y[1])
plt.axis("equal")
plt.show()
\end{verbatim}

\subsection*{Common warnings}

Arrow length in a plot may be automatically rescaled.

A normalized quiver plot hides speed.

A stream plot may hide sources, sinks, or singularities.

A vector field with a hole in its domain should not be plotted as if the hole were
ordinary.

% ---------------------------------------------------------------------
\section{Numerical differentiation and the Jacobian}
\label{app:G.3}

\subsection*{Forward difference}

For a one-variable function,
\[
\boxed{
    f'(a)\approx \frac{f(a+h)-f(a)}{h}.
}
\]

This is simple but less accurate.

\subsection*{Central difference}

\[
\boxed{
    f'(a)\approx
    \frac{f(a+h)-f(a-h)}{2h}.
}
\]

This is usually better.

\subsection*{Second derivative}

\[
\boxed{
    f''(a)\approx
    \frac{f(a+h)-2f(a)+f(a-h)}{h^2}.
}
\]

\subsection*{Partial derivatives}

For
\[
    f=f(x,y),
\]
\[
\boxed{
    f_x(a,b)
    \approx
    \frac{f(a+h,b)-f(a-h,b)}{2h}.
}
\]

\[
\boxed{
    f_y(a,b)
    \approx
    \frac{f(a,b+h)-f(a,b-h)}{2h}.
}
\]

\subsection*{Gradient by central differences}

\begin{verbatim}
def grad2(f, x, y, h=1e-5):
    fx = (f(x+h, y) - f(x-h, y)) / (2*h)
    fy = (f(x, y+h) - f(x, y-h)) / (2*h)
    return np.array([fx, fy])
\end{verbatim}

\subsection*{Jacobian matrix}

For
\[
    F:\mathbb R^n\to\mathbb R^m,
\]
the Jacobian matrix is
\[
\boxed{
    J_F(\mathbf p)
    =
    \begin{pmatrix}
    \dfrac{\partial F_1}{\partial x_1} & \cdots & \dfrac{\partial F_1}{\partial x_n}\\[2mm]
    \vdots & \ddots & \vdots\\[2mm]
    \dfrac{\partial F_m}{\partial x_1} & \cdots & \dfrac{\partial F_m}{\partial x_n}
    \end{pmatrix}_{\mathbf p}.
}
\]

\subsection*{Numerical Jacobian}

\begin{verbatim}
def jacobian(F, p, h=1e-5):
    p = np.array(p, dtype=float)
    m = len(F(p))
    n = len(p)
    J = np.zeros((m, n))

    for j in range(n):
        e = np.zeros(n)
        e[j] = 1.0
        J[:, j] = (F(p + h*e) - F(p - h*e)) / (2*h)

    return J
\end{verbatim}

\subsection*{Numerical Hessian}

For
\[
    f=f(x,y),
\]
\[
\boxed{
    f_{xx}\approx
    \frac{f(x+h,y)-2f(x,y)+f(x-h,y)}{h^2}.
}
\]

\[
\boxed{
    f_{yy}\approx
    \frac{f(x,y+h)-2f(x,y)+f(x,y-h)}{h^2}.
}
\]

\[
\boxed{
    f_{xy}\approx
    \frac{
    f(x+h,y+h)-f(x+h,y-h)-f(x-h,y+h)+f(x-h,y-h)
    }{4h^2}.
}
\]

\subsection*{Symbolic Jacobian}

Using \texttt{SymPy}:

\begin{verbatim}
import sympy as sp

x, y = sp.symbols("x y")

F = sp.Matrix([
    x**2 + y,
    sp.sin(x*y)
])

J = F.jacobian([x, y])
\end{verbatim}

\subsection*{Symbolic Hessian}

\begin{verbatim}
f = x**3 + x*y + y**2
H = sp.hessian(f, (x, y))
\end{verbatim}

\subsection*{Step-size warning}

If \(h\) is too large, the approximation error is large.

If \(h\) is too small, roundoff error can dominate.

A typical starting value is
\[
\boxed{
    h=10^{-5}
}
\]
for central differences, but this is not universal.

\subsection*{Differentiability warning}

Numerical derivatives can produce misleading numbers near:

\[
\begin{array}{ll}
1. & \text{corners,}\\
2. & \text{cusps,}\\
3. & \text{jumps,}\\
4. & \text{singularities,}\\
5. & \text{rapid oscillations.}
\end{array}
\]

A finite-difference approximation does not prove differentiability.

% ---------------------------------------------------------------------
\section{Numerical double and triple integration}
\label{app:G.4}

\subsection*{Midpoint rule in one variable}

\[
\boxed{
    \int_a^b f(x)\,dx
    \approx
    \sum_{i=1}^{N} f(x_i^*)\,\Delta x.
}
\]

Here
\[
    \Delta x=\frac{b-a}{N},
\]
and \(x_i^*\) is the midpoint of the \(i\)th subinterval.

\subsection*{Midpoint rule over a rectangle}

For
\[
    R=[a,b]\times[c,d],
\]
\[
\boxed{
    \iint_R f(x,y)\,dA
    \approx
    \sum_{i=1}^{M}\sum_{j=1}^{N}
    f(x_i^*,y_j^*)\,\Delta x\,\Delta y.
}
\]

\subsection*{Code for rectangular midpoint rule}

\begin{verbatim}
def midpoint_double(f, a, b, c, d, M=200, N=200):
    dx = (b - a) / M
    dy = (d - c) / N

    x = a + (np.arange(M) + 0.5)*dx
    y = c + (np.arange(N) + 0.5)*dy

    X, Y = np.meshgrid(x, y)
    return np.sum(f(X, Y))*dx*dy
\end{verbatim}

\subsection*{Region under variable bounds}

For
\[
    a\le x\le b,
    \qquad
    g_1(x)\le y\le g_2(x),
\]
\[
\boxed{
    \iint_R f\,dA
    =
    \int_a^b\int_{g_1(x)}^{g_2(x)}f(x,y)\,dy\,dx.
}
\]

Using \texttt{dblquad}:

\begin{verbatim}
from scipy.integrate import dblquad

value, error = dblquad(
    lambda y, x: f(x, y),
    a, b,
    lambda x: g1(x),
    lambda x: g2(x)
)
\end{verbatim}

Notice the order in the function:
\[
\boxed{
    \texttt{lambda y, x}
}
\]
because \texttt{dblquad} uses the inner variable first.

\subsection*{Polar numerical integration}

If
\[
    x=r\cos\theta,
    \qquad
    y=r\sin\theta,
\]
then
\[
\boxed{
    dA=r\,dr\,d\theta.
}
\]

So
\[
\boxed{
    \iint_R f(x,y)\,dA
    =
    \iint_{\widetilde R}
    f(r\cos\theta,r\sin\theta)\,r\,dr\,d\theta.
}
\]

Code pattern:

\begin{verbatim}
value, error = dblquad(
    lambda r, theta: f(r*np.cos(theta), r*np.sin(theta))*r,
    theta0, theta1,
    lambda theta: r0(theta),
    lambda theta: r1(theta)
)
\end{verbatim}

\subsection*{Triple integral over a box}

For
\[
    B=[a,b]\times[c,d]\times[e,f],
\]
\[
\boxed{
    \iiint_B F(x,y,z)\,dV
    \approx
    \sum F(x_i^*,y_j^*,z_k^*)\,\Delta x\,\Delta y\,\Delta z.
}
\]

\subsection*{Triple integral with variable bounds}

For
\[
    a\le x\le b,
\]
\[
    g_1(x)\le y\le g_2(x),
\]
\[
    h_1(x,y)\le z\le h_2(x,y),
\]
\[
\boxed{
    \iiint_E F\,dV
    =
    \int_a^b\int_{g_1(x)}^{g_2(x)}
    \int_{h_1(x,y)}^{h_2(x,y)}
    F(x,y,z)\,dz\,dy\,dx.
}
\]

Using \texttt{tplquad}:

\begin{verbatim}
from scipy.integrate import tplquad

value, error = tplquad(
    lambda z, y, x: F(x, y, z),
    a, b,
    lambda x: g1(x),
    lambda x: g2(x),
    lambda x, y: h1(x, y),
    lambda x, y: h2(x, y)
)
\end{verbatim}

Again, the innermost variable comes first.

\subsection*{Cylindrical numerical integration}

\[
\boxed{
    dV=r\,dr\,d\theta\,dz.
}
\]

\[
\boxed{
    F(x,y,z)
    =
    F(r\cos\theta,r\sin\theta,z).
}
\]

Use the integrand
\[
\boxed{
    F(r\cos\theta,r\sin\theta,z)\,r.
}
\]

\subsection*{Spherical numerical integration}

\[
\boxed{
    dV=\rho^2\sin\phi\,d\rho\,d\phi\,d\theta.
}
\]

Use the integrand
\[
\boxed{
    F(\rho\sin\phi\cos\theta,\rho\sin\phi\sin\theta,\rho\cos\phi)
    \rho^2\sin\phi.
}
\]

\subsection*{Monte Carlo estimate}

For a region \(R\) inside a box \(B\):

\[
\boxed{
    \iint_R f\,dA
    \approx
    \operatorname{Area}(B)
    \cdot
    \frac1N
    \sum_{i=1}^N f(x_i,y_i)\mathbf 1_R(x_i,y_i).
}
\]

Code pattern:

\begin{verbatim}
def monte_carlo_double(f, inside, a, b, c, d, N=100000):
    x = np.random.uniform(a, b, N)
    y = np.random.uniform(c, d, N)

    mask = inside(x, y)
    box_area = (b - a)*(d - c)

    return box_area*np.mean(f(x, y)*mask)
\end{verbatim}

Monte Carlo is usually slow but useful in irregular regions or high dimensions.

\subsection*{Common warnings}

Numerical integration returns an approximation and often an error estimate.

Singularities require special care.

Highly oscillatory functions require finer sampling.

Coordinate changes require Jacobian factors.

Improper integrals may appear finite on a finite sample but still diverge.

% ---------------------------------------------------------------------
\section{Symbolic line and surface integrals}
\label{app:G.5}

\subsection*{Symbolic setup}

\begin{verbatim}
import sympy as sp
\end{verbatim}

Define variables:
\[
\boxed{
    \texttt{x, y, z, t, u, v = sp.symbols("x y z t u v")}
}
\]

\subsection*{Line integral of a \texorpdfstring{\(1\)}{1}-form}

Let
\[
    \omega=P\,dx+Q\,dy+R\,dz.
\]

Let
\[
    \mathbf r(t)=(x(t),y(t),z(t)),
    \qquad
    a\le t\le b.
\]

Then
\[
\boxed{
    \int_C\omega
    =
    \int_a^b
    \left[
    P(\mathbf r(t))x'(t)
    +
    Q(\mathbf r(t))y'(t)
    +
    R(\mathbf r(t))z'(t)
    \right]\,dt.
}
\]

\subsection*{SymPy line integral pattern}

\begin{verbatim}
x, y, z, t = sp.symbols("x y z t")

P = ...
Q = ...
R = ...

xt = ...
yt = ...
zt = ...

integrand = (
    P.subs({x: xt, y: yt, z: zt})*sp.diff(xt, t)
    + Q.subs({x: xt, y: yt, z: zt})*sp.diff(yt, t)
    + R.subs({x: xt, y: yt, z: zt})*sp.diff(zt, t)
)

value = sp.integrate(integrand, (t, a, b))
\end{verbatim}

\subsection*{Line integral of a vector field}

For
\[
    \mathbf F=(P,Q,R),
\]
\[
\boxed{
    \int_C\mathbf F\cdot d\mathbf r
}
\]
uses the same formula as the \(1\)-form
\[
    P\,dx+Q\,dy+R\,dz.
\]

\subsection*{Surface flux integral}

Let
\[
    \mathbf F=(P,Q,R).
\]

Let
\[
    \mathbf r(u,v)=(x(u,v),y(u,v),z(u,v)).
\]

Then
\[
\boxed{
    \iint_S\mathbf F\cdot\widehat{\mathbf n}\,dS
    =
    \iint_D
    \mathbf F(\mathbf r(u,v))
    \cdot
    (\mathbf r_u\times\mathbf r_v)
    \,du\,dv.
}
\]

\subsection*{SymPy surface flux pattern}

\begin{verbatim}
x, y, z, u, v = sp.symbols("x y z u v")

P = ...
Q = ...
R = ...

xu = ...
yu = ...
zu = ...

r = sp.Matrix([xu, yu, zu])
ru = r.diff(u)
rv = r.diff(v)

normal = ru.cross(rv)

F = sp.Matrix([
    P.subs({x: xu, y: yu, z: zu}),
    Q.subs({x: xu, y: yu, z: zu}),
    R.subs({x: xu, y: yu, z: zu})
])

integrand = sp.simplify(F.dot(normal))

value = sp.integrate(
    integrand,
    (u, u0, u1),
    (v, v0, v1)
)
\end{verbatim}

The order of \((u,v)\) bounds must match the expression.

\subsection*{Surface area}

\[
\boxed{
    \operatorname{Area}(S)
    =
    \iint_D
    |\mathbf r_u\times\mathbf r_v|
    \,du\,dv.
}
\]

Symbolic pattern:

\begin{verbatim}
area_integrand = sp.sqrt(normal.dot(normal))
area = sp.integrate(area_integrand, (u, u0, u1), (v, v0, v1))
\end{verbatim}

\subsection*{Curl}

For
\[
    \mathbf F=(P,Q,R),
\]
\[
\boxed{
    \nabla\times\mathbf F
    =
    (R_y-Q_z,\ P_z-R_x,\ Q_x-P_y).
}
\]

SymPy pattern:

\begin{verbatim}
curlF = sp.Matrix([
    sp.diff(R, y) - sp.diff(Q, z),
    sp.diff(P, z) - sp.diff(R, x),
    sp.diff(Q, x) - sp.diff(P, y)
])
\end{verbatim}

\subsection*{Divergence}

\[
\boxed{
    \nabla\cdot\mathbf F=P_x+Q_y+R_z.
}
\]

SymPy pattern:

\begin{verbatim}
divF = sp.diff(P, x) + sp.diff(Q, y) + sp.diff(R, z)
\end{verbatim}

\subsection*{Gradient}

\[
\boxed{
    \nabla f=(f_x,f_y,f_z).
}
\]

SymPy pattern:

\begin{verbatim}
gradf = sp.Matrix([
    sp.diff(f, x),
    sp.diff(f, y),
    sp.diff(f, z)
])
\end{verbatim}

\subsection*{Jacobian}

For
\[
    F=(F_1,\ldots,F_m),
\]
\[
\boxed{
    J_F=\left(\frac{\partial F_i}{\partial x_j}\right).
}
\]

SymPy pattern:

\begin{verbatim}
F = sp.Matrix([F1, F2, F3])
J = F.jacobian([x, y, z])
\end{verbatim}

\subsection*{Green's theorem check}

For
\[
    \omega=P\,dx+Q\,dy,
\]
\[
\boxed{
    d\omega=(Q_x-P_y)\,dx\wedge dy.
}
\]

Symbolic scalar curl:

\begin{verbatim}
scalar_curl = sp.diff(Q, x) - sp.diff(P, y)
\end{verbatim}

\subsection*{Orientation warnings}

For a curve, reversing the parametrization changes the sign of a \(1\)-form integral.

For a surface, switching \(u\) and \(v\) changes
\[
    \mathbf r_u\times\mathbf r_v
\]
to
\[
    -\mathbf r_u\times\mathbf r_v.
\]

Flux changes sign when orientation reverses.

Surface area does not.

\subsection*{Symbolic warnings}

A symbolic answer may be correct but hard to read.

Use:
\[
\boxed{
    \texttt{sp.simplify}
}
\]

\[
\boxed{
    \texttt{sp.factor}
}
\]

\[
\boxed{
    \texttt{sp.expand}
}
\]

Always check assumptions, especially signs and parameter ranges.

% ---------------------------------------------------------------------
\section{Simulating flows and conservation laws}
\label{app:G.6}

\subsection*{Flow lines}

A flow line of a vector field
\[
    \mathbf F
\]
solves the differential equation
\[
\boxed{
    \frac{d\mathbf r}{dt}=\mathbf F(\mathbf r(t)).
}
\]

In the plane:
\[
\boxed{
    x'=P(x,y),
    \qquad
    y'=Q(x,y).
}
\]

In space:
\[
\boxed{
    x'=P(x,y,z),
    \qquad
    y'=Q(x,y,z),
    \qquad
    z'=R(x,y,z).
}
\]

\subsection*{Euler method}

For
\[
    \mathbf r'=\mathbf F(\mathbf r),
\]
Euler's method is
\[
\boxed{
    \mathbf r_{n+1}
    =
    \mathbf r_n+\Delta t\,\mathbf F(\mathbf r_n).
}
\]

It is simple but not very accurate.

\subsection*{Euler method code}

\begin{verbatim}
def euler_flow(F, r0, dt, steps):
    r = np.array(r0, dtype=float)
    points = [r.copy()]

    for _ in range(steps):
        r = r + dt*np.array(F(r))
        points.append(r.copy())

    return np.array(points)
\end{verbatim}

\subsection*{Using an ODE solver}

\begin{verbatim}
from scipy.integrate import solve_ivp

def ode(t, r):
    x, y = r
    return [Pfun(x, y), Qfun(x, y)]

sol = solve_ivp(
    ode,
    [t0, t1],
    [x0, y0],
    max_step=0.02
)

plt.plot(sol.y[0], sol.y[1])
plt.axis("equal")
plt.show()
\end{verbatim}

\subsection*{Several initial points}

\begin{verbatim}
for r0 in initial_points:
    sol = solve_ivp(ode, [t0, t1], r0, max_step=0.02)
    plt.plot(sol.y[0], sol.y[1])

plt.axis("equal")
plt.show()
\end{verbatim}

\subsection*{Equilibrium points}

Equilibrium points satisfy
\[
\boxed{
    \mathbf F(\mathbf p)=\mathbf 0.
}
\]

In the plane:
\[
\boxed{
    P(x,y)=0,
    \qquad
    Q(x,y)=0.
}
\]

Use symbolic solving:

\begin{verbatim}
solutions = sp.solve([P, Q], [x, y])
\end{verbatim}

Or numerical root finding if symbolic solving is too hard.

\subsection*{Linearization near an equilibrium}

If
\[
    \mathbf F(\mathbf p)=\mathbf 0,
\]
then near \(\mathbf p\),
\[
\boxed{
    \mathbf r'
    \approx
    J_{\mathbf F}(\mathbf p)(\mathbf r-\mathbf p).
}
\]

The eigenvalues of
\[
    J_{\mathbf F}(\mathbf p)
\]
help classify the local behavior.

\[
\begin{array}{c|c}
\text{Eigenvalue pattern} & \text{Typical behavior} \\ \hline
\text{both negative real parts} & \text{sink}\\[1mm]
\text{both positive real parts} & \text{source}\\[1mm]
\text{opposite signs} & \text{saddle}\\[1mm]
\text{pure imaginary} & \text{center or rotation-type behavior}\\[1mm]
\text{complex with negative real part} & \text{spiral sink}\\[1mm]
\text{complex with positive real part} & \text{spiral source}
\end{array}
\]

\subsection*{Conservation law}

A conservation law has the form
\[
\boxed{
    \frac{\partial \rho}{\partial t}
    +
    \nabla\cdot\mathbf J
    =
    s.
}
\]

Here:

\[
\begin{array}{c|c}
\rho & \text{density}\\[1mm]
\mathbf J & \text{flux density}\\[1mm]
s & \text{source rate}
\end{array}
\]

If there are no sources or sinks:
\[
\boxed{
    \rho_t+\nabla\cdot\mathbf J=0.
}
\]

\subsection*{Integral balance law}

For a fixed region \(E\),
\[
\boxed{
    \frac{d}{dt}\iiint_E \rho\,dV
    =
    -\iint_{\partial E}\mathbf J\cdot\widehat{\mathbf n}\,dS
    +
    \iiint_E s\,dV.
}
\]

In words:

\[
\boxed{
    \text{rate of stored amount}
    =
    -\text{outward flux}
    +
    \text{source production}.
}
\]

\subsection*{Numerical balance check}

At two nearby times,
\[
\boxed{
    \frac{A(t+\Delta t)-A(t)}{\Delta t}
}
\]
approximates
\[
    \frac{dA}{dt}.
\]

Compare it to
\[
\boxed{
    -\text{outward flux}+\text{source integral}.
}
\]

This is a useful diagnostic for simulations.

\subsection*{Discrete divergence in two dimensions}

On a grid with spacing \(\Delta x,\Delta y\), for
\[
    \mathbf J=(M,N),
\]
a central-difference approximation is
\[
\boxed{
    \nabla\cdot\mathbf J
    \approx
    \frac{M(x+\Delta x,y)-M(x-\Delta x,y)}{2\Delta x}
    +
    \frac{N(x,y+\Delta y)-N(x,y-\Delta y)}{2\Delta y}.
}
\]

\subsection*{Discrete divergence code}

\begin{verbatim}
def divergence2(M, N, dx, dy):
    M_y, M_x = np.gradient(M, dy, dx)
    N_y, N_x = np.gradient(N, dy, dx)
    return M_x + N_y
\end{verbatim}

Array-axis warning:
\[
\boxed{
    \texttt{np.gradient} \text{ usually returns } y\text{-derivatives first.}
}
\]

\subsection*{Advection by a velocity field}

If density moves with velocity
\[
    \mathbf v,
\]
then the flux density is
\[
\boxed{
    \mathbf J=\rho\mathbf v.
}
\]

With no sources:
\[
\boxed{
    \rho_t+\nabla\cdot(\rho\mathbf v)=0.
}
\]

If
\[
    \nabla\cdot\mathbf v=0
\]
and \(\rho\) is constant along moving particles, the flow is incompressible.

\subsection*{Common simulation warnings}

Euler's method can be unstable for large time steps.

A good-looking picture may violate conservation.

Numerical diffusion can smear sharp features.

Boundary conditions matter.

A vector field with a singularity should not be simulated through the singular point.

Always test a simulation on a case with a known answer.

% ---------------------------------------------------------------------
\section*{One-page computational summary}

\[
\boxed{
    \texttt{X, Y = np.meshgrid(x, y)}
}
\]

\[
\boxed{
    \texttt{ax.plot\_surface(X, Y, Z)}
}
\]

\[
\boxed{
    \texttt{plt.contour(X, Y, Z)}
}
\]

\[
\boxed{
    \texttt{plt.quiver(X, Y, P, Q)}
}
\]

\[
\boxed{
    \texttt{plt.streamplot(X, Y, P, Q)}
}
\]

\[
\boxed{
    f'(a)\approx\frac{f(a+h)-f(a-h)}{2h}
}
\]

\[
\boxed{
    f_x(a,b)\approx\frac{f(a+h,b)-f(a-h,b)}{2h}
}
\]

\[
\boxed{
    J_F(\mathbf p)_{ij}
    \approx
    \frac{F_i(\mathbf p+h\mathbf e_j)-F_i(\mathbf p-h\mathbf e_j)}{2h}
}
\]

\[
\boxed{
    \iint_R f\,dA
    \approx
    \sum f(x_i^*,y_j^*)\,\Delta x\,\Delta y
}
\]

\[
\boxed{
    dA=r\,dr\,d\theta
}
\]

\[
\boxed{
    dV=r\,dr\,d\theta\,dz
}
\]

\[
\boxed{
    dV=\rho^2\sin\phi\,d\rho\,d\phi\,d\theta
}
\]

\[
\boxed{
    \int_C P\,dx+Q\,dy+R\,dz
    =
    \int_a^b(Px'+Qy'+Rz')\,dt
}
\]

\[
\boxed{
    \iint_S\mathbf F\cdot\widehat{\mathbf n}\,dS
    =
    \iint_D
    \mathbf F(\mathbf r(u,v))\cdot(\mathbf r_u\times\mathbf r_v)
    \,du\,dv
}
\]

\[
\boxed{
    \mathbf r'=\mathbf F(\mathbf r)
}
\]

\[
\boxed{
    \rho_t+\nabla\cdot\mathbf J=s
}
\]

\[
\boxed{
    \frac{d}{dt}\int_E\rho\,dV
    =
    -\int_{\partial E}\mathbf J\cdot\widehat{\mathbf n}\,dS
    +
    \int_E s\,dV
}
\]